% Conclusion

We document that the securities lending market serves as an information acquisition channel for active mutual fund managers. Using unique regulatory data from the Deutsche Bundesbank with security-level lending positions, we show that active fund managers reduce their holdings in stocks they lend. The effect is economically large---a 0.33 percentage point position reduction per unit increase in lending---and is absent for passive funds, who cannot adjust their portfolios in response to the signal.

We identify the mechanism through mandatory short position disclosure events under EU regulation. Active funds reduce positions substantially more when lending precedes a disclosure event, confirming that the information content of lending drives the trading response. The timing is consistent with anticipation: only the lead disclosure amplifies the effect, while contemporaneous and lagged disclosures are uninformative. Results are robust to using passive fund lending as an exogenous source of variation in lending activity.

We further show that information flows within fund families, with the strongest effects from same-manager fund lending. The trading response is amplified in concentrated lending markets where individual transactions carry more informational weight.

Our findings contribute to the understanding of how information flows through financial markets. The securities lending market, typically studied from the perspective of short sellers, also functions as an information channel for the long side of the market. Fund managers who participate in lending gain private, real-time information about short-selling demand that enables them to make more informed portfolio decisions. These results have implications for the regulation of securities lending, the organization of fund families, and the broader debate about the role of financial market infrastructure in promoting price efficiency.
